%
% fig/argument.tex -- Geometrische Anschauung des Argumentprinzips
%
\begin{figure}
\centering
\includegraphics{papers/weyl/images/argument.pdf}
\caption{Geometrische Bedeutung des Wegintegrals: läuft $z$ einmal um den Rand
$\partial D$ eines Gebiets, so umrundet der Bildpunkt $f(z)$ den Ursprung in der
Bildebene genau so oft, wie $f$ Nullstellen in $D$ besitzt. Oben: $D$ enthält
eine Nullstelle, die Bildkurve $f(\partial D)$ macht eine Drehung um $0$.
Unten: $D$ enthält drei Nullstellen, die Bildkurve umläuft den Ursprung dreimal.
\label{weyl:fig:argument}}
\end{figure}
